\documentclass[11pt,letterpaper]{article}
\usepackage[T1]{fontenc}
\usepackage[utf8]{inputenc}
\usepackage{lmodern}
\usepackage[margin=1in,headheight=15pt]{geometry}
\usepackage{amsmath,amssymb,bm}
\usepackage{booktabs,array,longtable}
\usepackage{enumitem}
\usepackage{xcolor}
\usepackage{microtype}
\usepackage{fancyhdr}
\usepackage{hyperref}
\definecolor{ink}{HTML}{15354A}
\hypersetup{colorlinks=true,linkcolor=ink,citecolor=ink,urlcolor=ink,
  pdftitle={Horseshoe Linear Regression and Variable Selection},
  pdfauthor={Research notes},pdfsubject={Model, Gibbs sampler, and posterior decisions}}
\pagestyle{fancy}
\fancyhf{}
\fancyhead[L]{\small\textsc{Horseshoe linear regression}}
\fancyhead[R]{\small Review draft}
\fancyfoot[C]{\thepage}
\setlength{\parindent}{0pt}
\setlength{\parskip}{6pt plus 1pt minus 1pt}
\setlist{itemsep=3pt,topsep=4pt}
\setcounter{tocdepth}{2}
\numberwithin{equation}{section}
\allowdisplaybreaks[1]
\newcommand{\R}{\mathbb{R}}
\newcommand{\E}{\mathbb{E}}
\newcommand{\Prob}{\mathbb{P}}
\newcommand{\Var}{\operatorname{Var}}
\newcommand{\Cov}{\operatorname{Cov}}
\newcommand{\diag}{\operatorname{diag}}
\newcommand{\tr}{\operatorname{tr}}
\newcommand{\IG}{\operatorname{IG}}
\newcommand{\Ga}{\operatorname{Gamma}}
\newcommand{\HC}{\operatorname{C}^{+}}
\newcommand{\Normal}{\mathcal{N}}
\newcommand{\ind}{\mathbf{1}}
\newcommand{\given}{\mid}
\newcommand{\btheta}{\bm{\theta}}
\newcommand{\bbeta}{\bm{\beta}}
\newcommand{\blambda}{\bm{\lambda}}
\newcommand{\bnu}{\bm{\nu}}
\newcommand{\ones}{\mathbf{1}}
\newcommand{\trans}{^{\mathsf T}}
\newcommand{\norm}[1]{\left\lVert#1\right\rVert}
\newcommand{\key}[1]{\par\smallskip\textbf{#1}\par\smallskip}

\begin{document}
\hypersetup{pageanchor=false}
\begin{titlepage}
\vspace*{0.65in}
{\color{ink}\Large Bayesian Linear Regression and Variable Selection\par}
\vspace{0.45in}
{\color{ink}\Huge\bfseries Horseshoe Linear Regression\par}
\vspace{0.16in}
{\color{ink}\LARGE The model, the Gibbs sampler,\par and variable selection\par}
\vspace{0.55in}
{\large Detailed research notes\par}
{\large Review draft --- October 4, 2026\par}
\vspace{0.6in}
\hrule
\vspace{0.2in}
\textbf{Purpose.} Develop Gaussian linear regression with a horseshoe prior from
the likelihood to an implementable Gibbs sampler. Every full conditional is
derived, the distribution conventions are fixed explicitly, and posterior
shrinkage is separated from the decision to select variables.

\textbf{Main convention.} The coefficient prior is scaled by the noise variance:
\(\beta_j\given\sigma^2,\tau^2,\lambda_j^2\sim
\Normal(0,\sigma^2\tau^2\lambda_j^2)\). The intercept is separate and unpenalized.
The main derivation uses a proper inverse-gamma prior on \(\sigma^2\), with
common alternatives explained later.

\textbf{Scope.} These are mathematical and implementation notes, not an empirical
comparison. Worked quantities are analytic illustrations. No simulated data,
fitted model, convergence result, or benchmark is reported.
\vfill
\small Editable source: \texttt{horseshoe\_linear\_regression.tex}\par
The bibliography links to the original methodological sources.
\end{titlepage}
\pagenumbering{roman}
\hypersetup{pageanchor=true}
\begingroup
\small
\setlength{\parskip}{0pt}
\tableofcontents
\endgroup
\clearpage
\pagenumbering{arabic}

\section{The question the horseshoe answers}

Suppose a response may depend on many predictors, but we expect only a small
number of substantial effects. Ordinary least squares estimates each coefficient
without using that expectation. When predictors are numerous or correlated,
the estimates can be unstable; when \(p\geq n\), an unregularized coefficient
vector may not even be identifiable.

The horseshoe expresses sparsity through a continuous prior. A shared
\emph{global} scale encourages coefficients to be small, while a separate
\emph{local} scale lets each coefficient escape strong shrinkage when supported
by the likelihood. This global--local construction is the central idea of
Carvalho, Polson, and Scott \cite{carvalho2009,carvalho2010}.

There are three different tasks:
\begin{enumerate}
\item \textbf{Estimation:} describe plausible values of each coefficient and
      of the regression function.
\item \textbf{Prediction:} describe responses at new predictor values, including
      both parameter uncertainty and observation noise.
\item \textbf{Selection:} decide which variables to report or retain, using a
      stated rule or loss function.
\end{enumerate}
A horseshoe posterior directly addresses the first two. The third requires an
additional decision: the ordinary horseshoe does not assign positive probability
to any coefficient being exactly zero.

\subsection{Reading guide}

Sections \ref{sec:model}--\ref{sec:shrinkage} explain the model and its shrinkage
behavior. Sections \ref{sec:augmentation}--\ref{sec:sweep} give the augmented
posterior and derive the complete Gibbs sweep. Sections \ref{sec:linearalgebra}
and \ref{sec:hyperparameters} address implementation and prior calibration.
Sections \ref{sec:inference}--\ref{sec:diagnostics} cover inference, selection,
and verification. The appendices distinguish alternative parameterizations and
provide a compact reference sheet.

\subsection{Notation and distribution conventions}

Vectors are columns. In particular, \(y\in\R^n\), \(X\in\R^{n\times p}\),
and \(\bbeta\in\R^p\). The \(i\)th row of \(X\) is \(x_i\trans\), and
\(X_j\) denotes its \(j\)th column. The intercept is \(\alpha\in\R\).
The symbol \(\mid\cdot\) means conditioning on the data, fixed hyperparameters,
and all current unknowns except the variable being updated.

\begin{center}
\begin{tabular}{@{}p{0.20\textwidth}p{0.73\textwidth}@{}}
\toprule
Symbol & Meaning \\
\midrule
\(n,p\) & Number of observations and number of penalized slopes; the intercept is not counted in \(p\). \\
\(\sigma^2\) & Observation-noise variance; \(\sigma=\sqrt{\sigma^2}>0\). \\
\(\tau,s_\tau\) & Global shrinkage scale and the fixed scale of its half-Cauchy prior. \\
\(\lambda_j\) & Local shrinkage scale for slope \(j\). \\
\(\nu_j,\xi\) & Positive auxiliary variables used to represent half-Cauchy priors. \\
\(a_\sigma,b_\sigma\) & Shape and inverse-gamma scale for the noise-variance prior. \\
\(D,A,m\) & Prior variance multiplier, conditional precision multiplier, and conditional coefficient mean. \\
\bottomrule
\end{tabular}
\end{center}

We use \(\Normal_k(m,V)\) with \emph{covariance} \(V\), not precision. The
inverse-gamma convention is
\begin{equation}\label{eq:ig}
 Z\sim\IG(a,b)
 \quad\Longleftrightarrow\quad
 p(z)=\frac{b^a}{\Gamma(a)}z^{-a-1}\exp(-b/z),\qquad z>0.
\end{equation}
Thus \(1/Z\sim\Ga(a,\text{rate}=b)\). If software uses gamma shape and scale,
generate \(G\sim\Ga(a,\text{scale}=1/b)\) and return \(1/G\).
The mean \(b/(a-1)\) exists only when \(a>1\); the mode is \(b/(a+1)\).

A half-Cauchy distribution here is parameterized by its \emph{scale} \(s\):
\begin{equation}\label{eq:hc}
 T\sim\HC(0,s),\qquad
 p(t)=\frac{2s}{\pi(s^2+t^2)},\qquad t>0.
\end{equation}
The symbols \(\tau\) and \(\lambda_j\) are standard-deviation multipliers;
\(\tau^2\) and \(\lambda_j^2\) are the variables sampled in the augmented Gibbs
scheme. Confusing a scale with its square changes the model.

\section{The regression model and the prior}\label{sec:model}

\subsection{Likelihood}

Assume conditionally independent, homoscedastic Gaussian errors:
\begin{align}
 y_i &= \alpha+x_i\trans\bbeta+\varepsilon_i,
 & \varepsilon_i\given\sigma^2&\overset{\mathrm{iid}}\sim\Normal(0,\sigma^2),\\
 y\given\alpha,\bbeta,\sigma^2
 &\sim\Normal_n(\alpha\ones_n+X\bbeta,\sigma^2 I_n).
\end{align}
Consequently,
\begin{equation}\label{eq:likelihood}
 p(y\given\alpha,\bbeta,\sigma^2)
 =(2\pi\sigma^2)^{-n/2}
 \exp\left\{-\frac{\norm{y-\alpha\ones_n-X\bbeta}^2}{2\sigma^2}\right\}.
\end{equation}
This likelihood is an assumption about the errors, not a requirement that
predictors themselves be normally distributed.

\subsection{The noise-scaled horseshoe hierarchy}

Our main model is
\begin{align}
 \beta_j\given\sigma^2,\tau^2,\lambda_j^2
 &\overset{\mathrm{ind}}\sim\Normal(0,\sigma^2\tau^2\lambda_j^2),
 &&j=1,\ldots,p,\label{eq:coefprior}\\
 \lambda_j&\overset{\mathrm{iid}}\sim\HC(0,1),
 &\tau&\sim\HC(0,s_\tau),\label{eq:scalepriors}\\
 \sigma^2&\sim\IG(a_\sigma,b_\sigma),
 &p(\alpha)&\propto1,\label{eq:otherpriors}
\end{align}
where \(a_\sigma,b_\sigma,s_\tau>0\). The scale priors and the noise-variance
prior are mutually independent. The slopes are independent \emph{conditional}
on their scales and \(\sigma^2\); they are not generally independent after
integrating out the shared scales.

Define
\begin{equation}\label{eq:D}
 D=\tau^2\diag(\lambda_1^2,\ldots,\lambda_p^2),
 \qquad \bbeta\given\sigma^2,D\sim\Normal_p(0,\sigma^2D).
\end{equation}
The prior precision for coefficient \(j\) is
\(1/(\sigma^2\tau^2\lambda_j^2)\). Small \(\tau\) encourages shrinkage
throughout the model. A large \(\lambda_j\) weakens the penalty for its own
coefficient. A small \(\lambda_j\) strengthens that penalty.

The conventional unit global scale is \(s_\tau=1\). It is a transparent
reference choice, not a universal calibration. Section \ref{sec:hyperparameters}
explains how sample size, predictor scaling, and an anticipated degree of
sparsity affect a more informed choice.

\key{A notation boundary for this project}
The \(\lambda_j\)'s in these notes are horseshoe local scales with half-Cauchy
priors. They are not the lambda parameter of a quantile-regression sampler.
The project's Gamma\((1,\text{rate}=\sqrt n)\) convention for that different
parameter does not replace the horseshoe hierarchy.

\subsection{Centering, scaling, and the intercept}

Predictor units affect the meaning of a prior on coefficients. Unless units
carry a reason for unequal prior scales, center each nonconstant predictor and
divide it by a recorded scale. One precise convention is
\begin{equation}\label{eq:standardize}
 \bar x_j=\frac1n\sum_i x_{ij}^{\rm raw},\qquad
 s_j^2=\frac1n\sum_i(x_{ij}^{\rm raw}-\bar x_j)^2,\qquad
 x_{ij}=\frac{x_{ij}^{\rm raw}-\bar x_j}{s_j}.
\end{equation}
Then \(\ones_n\trans X=0\) and \(X_j\trans X_j=n\). Using sample standard
deviation instead gives \(X_j\trans X_j=n-1\); either is valid if used
consistently. Constant columns must be removed or handled separately.

An intercept is an overall baseline, not a candidate sparse predictor. We
therefore keep it outside \(D\). Under centered predictors,
\(\alpha\given\cdot\sim\Normal(\bar y,\sigma^2/n)\). Keeping this draw
retains baseline uncertainty for prediction.

Centering \(y\) and setting its intercept to zero is not automatically
equivalent to integrating out an unknown intercept. With a flat intercept
prior, exact integration reduces the residual dimension from \(n\) to
\(n-1\). Appendix \ref{app:intercept} gives the resulting formulas. The main
sampler avoids ambiguity by sampling \(\alpha\) explicitly.

For validation or prediction, estimate preprocessing constants using training
data only, and apply those same constants to held-out observations. If the
response was also transformed as \(y'=(y^{\rm raw}-c_y)/s_y\), then
\begin{equation}\label{eq:backtransform}
 \beta_j^{\rm raw}=\frac{s_y}{s_j}\beta_j,\qquad
 \alpha^{\rm raw}=c_y+s_y\alpha-
       \sum_j\frac{s_y\beta_j\bar x_j}{s_j},\qquad
 \sigma^{\rm raw}=s_y\sigma.
\end{equation}
Apply these transformations to every posterior draw, so uncertainty transforms
with the parameters.

\subsection{What makes this prior a horseshoe?}

Conditional normality makes computation tractable. Mixing that normal over a
half-Cauchy local scale gives a marginal coefficient density with an unbounded
spike at zero and heavy tails. For fixed positive \(\tau,\sigma\), it remains
a proper continuous density. An infinite density value at one point is not a
positive probability at that point.

The spike supports estimates close to zero; the tails accommodate large
effects without a Gaussian prior's strong tail penalty. This does not guarantee
that a weak signal is retained or that a particular correlated predictor is
identified. The likelihood and the global-scale prior still matter.

\section{Understanding shrinkage through an orthogonal design}\label{sec:shrinkage}

\subsection{A one-coordinate normal calculation}

For intuition, suppose the predictor columns are mutually orthogonal and
orthogonal to the intercept. Let \(c_j=X_j\trans X_j>0\). The least-squares
coordinate is
\[
 \widehat\beta_j=\frac{X_j\trans(y-\alpha\ones_n)}{c_j},
 \qquad
 \widehat\beta_j\given\beta_j,\sigma^2\sim
 \Normal\left(\beta_j,\frac{\sigma^2}{c_j}\right).
\]
Combining this normal likelihood with \eqref{eq:coefprior} adds their
precisions. Define the shrinkage factor
\begin{equation}\label{eq:kappa}
 \kappa_j=\frac{1}{1+c_j\tau^2\lambda_j^2}.
\end{equation}
The conditional posterior becomes
\begin{equation}\label{eq:orthposterior}
 \beta_j\given\cdot\sim
 \Normal\left((1-\kappa_j)\widehat\beta_j,
      \frac{\sigma^2}{c_j}(1-\kappa_j)\right).
\end{equation}
Thus \(\kappa_j\) near one means strong shrinkage; \(\kappa_j\) near zero
means weak shrinkage. The noise variance cancels out of \(\kappa_j\) because
the likelihood and coefficient prior share the factor \(\sigma^2\).

For a unit-length column, \(c_j=1\). Under \eqref{eq:standardize}, \(c_j=n\).
Using \(1/(1+\tau^2\lambda_j^2)\) indiscriminately in standardized regression
misses this factor of \(n\).

\subsection{Why the name? A transformation of the local scale}

For a unit-information coordinate with \(c_j=1\) and fixed \(\tau=1\),
\(\lambda_j\sim\HC(0,1)\) implies
\(\kappa_j=1/(1+\lambda_j^2)\sim\operatorname{Beta}(1/2,1/2)\).
Its density is high near both endpoints. One end represents substantial
shrinkage and the other little shrinkage.

More generally put \(t=\sqrt{c_j}\tau\). Inverting
\(\kappa=1/(1+t^2\lambda^2)\) and including the Jacobian gives
\begin{equation}\label{eq:kappadensity}
 p(\kappa\given\tau)=
 \frac{t}{\pi\sqrt{\kappa(1-\kappa)}\{1+(t^2-1)\kappa\}},
 \qquad 0<\kappa<1.
\end{equation}
This is a \emph{prior} density at fixed \(\tau\). Data change the posterior
distribution of \(\kappa\); it is not a prior inclusion probability.

\subsection{An analytic illustration}

Take an orthogonal column with \(c_j=100\), and hold \(\tau=0.05\) fixed.
Then \(c_j\tau^2=0.25\). The following values are direct substitutions in
\eqref{eq:kappa}; they are not fitted estimates or simulation results.
\begin{center}
\begin{tabular}{@{}rrrr@{}}
\toprule
Local scale \(\lambda_j\) & \(c_j\tau^2\lambda_j^2\) & \(\kappa_j\) & Fraction retained \(1-\kappa_j\)\\
\midrule
0.2 & 0.01 & 0.9901 & 0.0099\\
1   & 0.25 & 0.8000 & 0.2000\\
10  & 25   & 0.0385 & 0.9615\\
\bottomrule
\end{tabular}
\end{center}
The same global scale can nearly eliminate one coordinate and preserve most
of another. In actual inference the scales are random and must be averaged
over, rather than fixed at convenient values.

\subsection{The limitation with correlated predictors}

For a general design, the conditional mean is
\((X\trans X+D^{-1})^{-1}X\trans(y-\alpha\ones_n)\). Off-diagonal entries of
\(X\trans X\) couple the coefficients. There is generally no exact formula
that independently multiplies each ordinary least-squares coefficient by
\(1-\kappa_j\).

In particular, two nearly interchangeable predictors can trade explanatory
roles across posterior draws. Their joint contribution may be stable while
their individual coefficients remain uncertain. Shrinkage factors remain
useful scale summaries, but their coordinatewise interpretation is exact only
in the orthogonal setting.

\section{Turning half-Cauchy priors into conjugate updates}\label{sec:augmentation}

\subsection{The auxiliary-variable identity}

The Gibbs sampler uses the inverse-gamma representation employed by Makalic
and Schmidt \cite{makalic2016}. For any \(s>0\), introduce \(V>0\) and write
\begin{equation}\label{eq:identity}
 Z\given V\sim\IG(1/2,1/V),\qquad
 V\sim\IG(1/2,1/s^2).
\end{equation}
Then \(\sqrt Z\sim\HC(0,s)\).

To verify the identity, retain all factors that depend on \(z\) or \(v\):
\[
 p(z,v)=\frac{1}{\pi s}\,z^{-3/2}v^{-2}
       \exp\left\{-\frac{1/z+1/s^2}{v}\right\}.
\]
For \(C>0\), substitution \(u=C/v\) gives
\(\int_0^\infty v^{-2}e^{-C/v}\,dv=1/C\). Hence
\begin{equation}
 p(z)=\frac{1}{\pi s}\frac{z^{-3/2}}{1/z+1/s^2}
     =\frac{s}{\pi\sqrt z(s^2+z)}.
\end{equation}
Setting \(z=t^2\), the Jacobian is \(2t\), so
\(p(t)=2s/\{\pi(s^2+t^2)\}\), exactly \eqref{eq:hc}.

The auxiliary variable has not approximated or replaced the half-Cauchy.
Integrating it out recovers the original prior exactly. Its purpose is to make
the conditional distributions easy to sample.

\subsection{The augmented horseshoe model}

Apply \eqref{eq:identity} independently to every local scale and once to the
global scale:
\begin{align}
 \lambda_j^2\given\nu_j&\sim\IG(1/2,1/\nu_j),
 &\nu_j&\sim\IG(1/2,1),\label{eq:localaug}\\
 \tau^2\given\xi&\sim\IG(1/2,1/\xi),
 &\xi&\sim\IG(1/2,1/s_\tau^2).\label{eq:globalaug}
\end{align}
Together with the Gaussian likelihood, coefficient prior, and
\eqref{eq:otherpriors}, this is the model sampled below. In this chosen
parameterization, \(s_\tau\) appears in the prior on \(\xi\), and later in
the full conditional for \(\xi\). Moving it into a different update without
redefining the augmentation would be an error.

The variables \(\nu_j\) and \(\xi\) are computational devices. They are
neither regression coefficients nor indicators that a predictor is included.
They can be omitted from scientific summaries, but must be maintained in the
chain state if these conditional updates are used.

\section{The posterior kernel and a recipe for deriving conditionals}\label{sec:kernel}

Let \(r=y-\alpha\ones_n-X\bbeta\). Up to constants independent of all
unknowns, the augmented joint posterior density with respect to
\(d\alpha\,d\bbeta\,d\sigma^2\,d\tau^2\,d\blambda^2\,d\bnu\,d\xi\) is
\begin{align}\label{eq:jointkernel}
 \pi(\alpha,\bbeta,\sigma^2,\tau^2,\blambda^2,\bnu,\xi\given y,X)
 &\propto
 (\sigma^2)^{-n/2}\exp\left(-\frac{r\trans r}{2\sigma^2}\right)\nonumber\\
 &\quad\times(\sigma^2)^{-p/2}(\tau^2)^{-p/2}
 \prod_{j=1}^p(\lambda_j^2)^{-1/2}
 \exp\left(-\sum_j\frac{\beta_j^2}{2\sigma^2\tau^2\lambda_j^2}\right)\nonumber\\
 &\quad\times(\sigma^2)^{-a_\sigma-1}\exp(-b_\sigma/\sigma^2)\nonumber\\
 &\quad\times\prod_{j=1}^p
 \left[\nu_j^{-1/2}(\lambda_j^2)^{-3/2}
       \exp\left(-\frac{1}{\nu_j\lambda_j^2}\right)
       \nu_j^{-3/2}\exp(-1/\nu_j)\right]\nonumber\\
 &\quad\times\xi^{-1/2}(\tau^2)^{-3/2}
       \exp\left(-\frac{1}{\xi\tau^2}\right)
       \xi^{-3/2}\exp\left(-\frac{1}{s_\tau^2\xi}\right).
\end{align}
The flat intercept prior contributes no factor. Fixed factors such as
\((s_\tau^2)^{-1/2}\) can be suppressed here.

To derive a full conditional, keep only factors involving the variable being
updated, holding all other variables fixed. For inverse-gamma updates,
identify the exponent of the variable and the coefficient of its reciprocal.
For normal updates, complete a quadratic square.

For example, the factor \(\nu_j^{-1/2}\) comes from the normalization of
\(\lambda_j^2\given\nu_j\). It cannot be discarded when updating \(\nu_j\).
Likewise, \((\sigma^2)^{-p/2}\) in the coefficient prior is essential when
updating \(\sigma^2\). Many incorrect samplers arise from treating a
normalizing constant as constant in the wrong variable.

\section{Derivation of every full conditional}\label{sec:conditionals}

\subsection{Intercept: one unpenalized normal draw}\label{sec:alpha}

Write \(z=y-X\bbeta\) and \(\bar z=n^{-1}\ones_n\trans z\). Then
\[
 \sum_i(z_i-\alpha)^2
 =\sum_i(z_i-\bar z)^2+n(\alpha-\bar z)^2.
\]
Only the last term involves \(\alpha\). With its flat prior,
\begin{equation}\label{eq:alpha}
 \boxed{\alpha\given\cdot\sim\Normal\left(
       \frac{\ones_n\trans(y-X\bbeta)}{n},\frac{\sigma^2}{n}\right).}
\end{equation}
For centered \(X\), the mean is simply \(\bar y\). If instead
\(\alpha\sim\Normal(m_\alpha,v_\alpha)\), independent of \(\sigma^2\), replace
the variance by \(V_\alpha=(n/\sigma^2+1/v_\alpha)^{-1}\) and the mean by
\(V_\alpha\{\ones_n\trans(y-X\bbeta)/\sigma^2+m_\alpha/v_\alpha\}\).
This independent proper intercept prior does not add a power of \(\sigma^2\).

\subsection{Regression coefficients: a multivariate normal block}

Set \(\widetilde y=y-\alpha\ones_n\). The terms involving \(\bbeta\) are
\[
 \pi(\bbeta\given\cdot)\propto
 \exp\left\{-\frac{1}{2\sigma^2}
   \left[(\widetilde y-X\bbeta)\trans(\widetilde y-X\bbeta)
          +\bbeta\trans D^{-1}\bbeta\right]\right\}.
\]
Expand the quadratic and define
\begin{equation}\label{eq:A}
 A=X\trans X+D^{-1},\qquad b=X\trans\widetilde y,\qquad m=A^{-1}b.
\end{equation}
Because \(Am=b\), completing the square gives
\[
 \bbeta\trans A\bbeta-2\bbeta\trans b
 =(\bbeta-m)\trans A(\bbeta-m)-m\trans Am.
\]
The final term is constant in \(\bbeta\), and therefore
\begin{equation}\label{eq:beta}
 \boxed{\bbeta\given\cdot\sim\Normal_p(m,\sigma^2 A^{-1}).}
\end{equation}
For any nonzero \(v\),
\(v\trans Av=\norm{Xv}^2+\sum_jv_j^2/(\tau^2\lambda_j^2)>0\).
Thus \(A\) is positive definite for finite positive scales, even if
\(X\trans X\) is singular. The conditional normal distribution is well
defined for \(p>n\). This is a statement about regularization, not a claim
that the data identify every coefficient separately.

Drawing the entire vector as a block respects its conditional dependence.
Replacing the draw by \(m\) would discard posterior variability and cease to
be a Gibbs sampler. Section \ref{sec:linearalgebra} describes how to draw this
normal without computing a matrix inverse.

\subsection{Noise variance: likelihood plus coefficient-prior contribution}

Define the two nonnegative sums
\begin{equation}\label{eq:RSSQ}
 R=\norm{y-\alpha\ones_n-X\bbeta}^2,\qquad
 Q=\bbeta\trans D^{-1}\bbeta
   =\sum_{j=1}^p\frac{\beta_j^2}{\tau^2\lambda_j^2}.
\end{equation}
Combining the likelihood, coefficient-prior normalization, and noise prior,
\[
 \pi(\sigma^2\given\cdot)\propto
 (\sigma^2)^{-\{a_\sigma+(n+p)/2\}-1}
 \exp\left\{-\frac{b_\sigma+(R+Q)/2}{\sigma^2}\right\}.
\]
Comparison with \eqref{eq:ig} yields
\begin{equation}\label{eq:sigma}
 \boxed{\sigma^2\given\cdot\sim
 \IG\left(a_\sigma+\frac{n+p}{2},\ b_\sigma+\frac{R+Q}{2}\right).}
\end{equation}
The \(p/2\) contribution is required because the prior variance of each of
the \(p\) slopes contains \(\sigma^2\). The term \(Q\) is required for the
same reason. The intercept contributes neither an extra \(1/2\) nor an
intercept penalty because its prior is flat. The likelihood still contributes
\(n/2\), since \(\alpha\) is conditioned on here, not integrated out.

\subsection{Local variance multiplier for each coefficient}

For \(\ell_j=\lambda_j^2\), the coefficient prior contributes
\(\ell_j^{-1/2}\exp\{-\beta_j^2/(2\sigma^2\tau^2\ell_j)\}\), and the
augmented prior contributes
\(\ell_j^{-3/2}\exp\{-1/(\nu_j\ell_j)\}\). Hence
\[
 \pi(\ell_j\given\cdot)\propto\ell_j^{-2}
 \exp\left\{-\frac{\nu_j^{-1}+\beta_j^2/(2\sigma^2\tau^2)}{\ell_j}\right\},
\]
which is
\begin{equation}\label{eq:lambda}
 \boxed{\lambda_j^2\given\cdot\sim
 \IG\left(1,\ \frac1{\nu_j}+\frac{\beta_j^2}{2\sigma^2\tau^2}\right),
 \qquad j=1,\ldots,p.}
\end{equation}
These updates are conditionally independent across \(j\). A large
\(\beta_j^2/(\sigma^2\tau^2)\) increases the inverse-gamma scale and allows
a larger local variance. This is how a coefficient can resist global
shrinkage. It is a stochastic update, not a deterministic monotone rule.

The conditional shape is one. Its mean is infinite, so an argument based on
the ``conditional mean of \(\lambda_j^2\)'' is not available. Simulating this
proper distribution is nevertheless well defined. Scale medians, log-scales,
and shrinkage factors are often more interpretable summaries.

\subsection{Global variance multiplier}

Put \(t=\tau^2\). All \(p\) coefficient priors contribute a total factor
\(t^{-p/2}\), while the augmented global prior contributes \(t^{-3/2}\).
Thus
\[
 \pi(t\given\cdot)\propto
 t^{-(p+3)/2}
 \exp\left\{-\frac{\xi^{-1}+
          \sum_j\beta_j^2/(2\sigma^2\lambda_j^2)}{t}\right\}.
\]
Matching the power \(-a-1\) gives
\begin{equation}\label{eq:tau}
 \boxed{\tau^2\given\cdot\sim
 \IG\left(\frac{p+1}{2},\ \frac1\xi+
           \frac1{2\sigma^2}\sum_{j=1}^p\frac{\beta_j^2}{\lambda_j^2}\right).}
\end{equation}
Unlike a local update, this update aggregates information from all slopes.
Large local scales can accommodate a few substantial coefficients while
allowing the global scale to remain small for the others.

\subsection{Auxiliary variables}

For \(\nu_j\), the normalization in
\(\lambda_j^2\given\nu_j\) supplies \(\nu_j^{-1/2}\). Multiplying by
its own prior gives
\[
 \pi(\nu_j\given\cdot)\propto
 \nu_j^{-2}\exp\left\{-\frac{1+1/\lambda_j^2}{\nu_j}\right\}.
\]
Similarly,
\[
 \pi(\xi\given\cdot)\propto
 \xi^{-2}\exp\left\{-\frac{1/s_\tau^2+1/\tau^2}{\xi}\right\}.
\]
Therefore,
\begin{align}
 \boxed{\nu_j\given\cdot\sim\IG\left(1,\ 1+\frac1{\lambda_j^2}\right),}
 &\qquad j=1,\ldots,p,\label{eq:nu}\\
 \boxed{\xi\given\cdot\sim\IG\left(1,\ \frac1{s_\tau^2}+\frac1{\tau^2}\right).}
 &\label{eq:xi}
\end{align}
When \(s_\tau=1\), the last scale reduces to \(1+1/\tau^2\).
No Metropolis acceptance step and no slice-sampling subroutine is needed for
any of these updates in the stated Gaussian model.

\section{How the Gibbs sampler works}\label{sec:sweep}

\subsection{From conditional draws to a joint posterior sample}

A Gibbs sampler repeatedly replaces one block of unknowns by a draw from its
conditional posterior given the current remaining blocks. If the chain is
already distributed according to the joint posterior, any one exact
conditional update leaves that joint distribution unchanged. A sequence of
such updates also leaves it unchanged.

This invariance identifies the target distribution; it does not say that a
chain started arbitrarily is immediately at stationarity. Nor does it make
successive draws independent. Appropriate regularity and adequate exploration
are needed for the empirical distribution of a long run to approximate the
posterior. Diagnostics assess finite-run behavior, not just the correctness
of the formulas.

\subsection{A complete sweep with explicit update order}

Store
\(\{\alpha,\bbeta,\sigma^2,\blambda^2,\tau^2,\bnu,\xi\}\).
Choose finite initial values with every variance and auxiliary variable
strictly positive. For example, start \(\alpha\) near \(\bar y\),
\(\bbeta\) near zero, \(\lambda_j^2=\nu_j=1\), a positive noise variance,
\(\tau^2=s_\tau^2\), and \(\xi=1\). These are starting values, not prior
draws or estimated hyperparameters. Use dispersed starts across chains.

At iteration \(t=1,\ldots,T\), carry out:
\begin{enumerate}[label=\textbf{\arabic*.},leftmargin=*]
\item \textbf{Update the intercept.} Draw \(\alpha^{(t)}\) using
\eqref{eq:alpha} with \(\bbeta^{(t-1)},\sigma^{2(t-1)}\).

\item \textbf{Update all slopes together.} Form \(D\) from
\(\tau^{2(t-1)},\blambda^{2(t-1)}\), set
\(\widetilde y=y-\alpha^{(t)}\ones_n\), and draw \(\bbeta^{(t)}\)
from \eqref{eq:beta} with \(\sigma^{2(t-1)}\).

\item \textbf{Update the noise variance.} Evaluate \(R\) using the new
\(\alpha^{(t)},\bbeta^{(t)}\), and \(Q\) using the still-current old scales.
Draw \(\sigma^{2(t)}\) from \eqref{eq:sigma}.

\item \textbf{Update local variances.} For every \(j\), draw
\(\lambda_j^{2(t)}\) from \eqref{eq:lambda} with
\(\beta_j^{(t)},\sigma^{2(t)},\tau^{2(t-1)},\nu_j^{(t-1)}\).
The local draws can be generated as one conditionally independent batch.

\item \textbf{Update the global variance.} Draw \(\tau^{2(t)}\) from
\eqref{eq:tau} using \(\bbeta^{(t)},\sigma^{2(t)},\blambda^{2(t)}\)
and \(\xi^{(t-1)}\).

\item \textbf{Update the auxiliary variables.} Draw \(\nu_j^{(t)}\)
from \eqref{eq:nu} using \(\lambda_j^{2(t)}\), and draw \(\xi^{(t)}\)
from \eqref{eq:xi} using \(\tau^{2(t)}\).

\item \textbf{Save the completed state.} After the chosen initial transient,
store draws and any desired predictive or shrinkage summaries. Save
checkpoints separately if the run must be restartable.
\end{enumerate}
Each update uses the most recently available values. Other sequential scan
orders can be valid, but updating every block simultaneously using only the
previous iteration's values is generally not the same Gibbs transition and
need not preserve the posterior.

\subsection{Why the updates adapt to sparsity}

The coefficient draw combines data fit with the present penalties. The local
updates then accommodate the magnitudes of individual coefficients. The
global update adjusts the overall scale after accounting for those local
exceptions. Finally, the auxiliary updates keep the scale-mixture
representation consistent with the intended half-Cauchy priors.

This feedback repeats throughout the run. One iteration is not a fitted
answer; the posterior distribution is represented by the retained collection
of complete states. A small coefficient at one iteration does not mean that
the variable has been permanently removed.

\subsection{An implementation skeleton}

The following is language-neutral pseudocode. \texttt{IGdraw(a,b)} follows
\eqref{eq:ig}, and \texttt{NormalBlock} performs the random draw in
\eqref{eq:beta}, not just its conditional mean.
\begin{quote}\small
\begin{verbatim}
precompute XtX = transpose(X) * X, Xty = transpose(X) * y
precompute Xt1 = transpose(X) * ones(n)
initialize alpha, beta, sigma2, lambda2, tau2, nu, xi
for iteration = 1,...,T:
    alpha = Normal(mean(y - X * beta), sigma2 / n)
    d = tau2 * lambda2
    A = XtX + diag(1 / d)
    b = Xty - alpha * Xt1
    beta = NormalBlock(A, b, sigma2)
    residual = y - alpha * ones(n) - X * beta
    R = dot(residual, residual)
    Q = sum(beta^2 / d)
    sigma2 = IGdraw(a_sigma + (n+p)/2, b_sigma + (R+Q)/2)
    lambda2 = IGdraw(1, 1/nu + beta^2/(2*sigma2*tau2))
    tau2 = IGdraw((p+1)/2, 1/xi + sum(beta^2/lambda2)/(2*sigma2))
    nu = IGdraw(1, 1 + 1/lambda2)
    xi = IGdraw(1, 1/s_tau^2 + 1/tau2)
    save completed state if retained
\end{verbatim}
\end{quote}
Arithmetic on vectors is elementwise except for explicit matrix products.
The second argument of \texttt{Normal} above is a variance. For \(p\gg n\),
the next section replaces the \(p\times p\) block construction, so
\(X\trans X\) need not be stored.

\section{Sampling the Gaussian block accurately and efficiently}\label{sec:linearalgebra}

\subsection{Cholesky sampling when the coefficient dimension is manageable}

Although \eqref{eq:beta} is written using \(A^{-1}\), implementation should
use linear solves. Let \(A=LL\trans\), where \(L\) is lower triangular.
First solve
\[
 Lq=b,\qquad L\trans m=q.
\]
Draw \(z\sim\Normal_p(0,I_p)\) and solve \(L\trans v=z\). Return
\begin{equation}\label{eq:choldraw}
 \bbeta=m+\sigma v.
\end{equation}
Indeed,
\[
 \Cov(v)=L^{-\mathsf T}I_pL^{-1}=(LL\trans)^{-1}=A^{-1}.
\]
The transpose matters. Solving \(Lv=z\) instead would generally give
\(L^{-1}L^{-\mathsf T}\), which is not \(A^{-1}\). If a software package
returns an upper-triangular factor \(U\) with \(A=U\trans U\), solve
\(Uv=z\) instead.

For a dense design, precomputing \(X\trans X\) costs \(O(np^2)\); the
factorization then costs \(O(p^3)\) per sweep, with \(O(p^2)\) storage for
the precision matrix. The changing scales require updating its diagonal and
refactorizing. The factorization can also supply a log determinant if needed:
\(\log|A|=2\sum_j\log L_{jj}\).

\subsection{An exact observation-space draw when \texorpdfstring{\(p\)}{p} is large}

Bhattacharya, Chakraborty, and Mallick \cite{bhattacharya2016} give an exact
Gaussian sampling method that can exploit \(n\ll p\). Here is its form for
our parameterization. Condition on \(\alpha,\sigma^2,D\) and put
\(H=I_n+XDX\trans\).
\begin{enumerate}
\item Draw independently \(u\sim\Normal_p(0,D)\) and
      \(\delta\sim\Normal_n(0,I_n)\).
\item Form \(v=Xu+\delta\).
\item Solve \(Hw=\widetilde y/\sigma-v\).
\item Return \(\bbeta=\sigma(u+DX\trans w)\).
\end{enumerate}
All normal variables in this construction are freshly drawn each sweep.
With diagonal \(D\), generate \(u_j=\sqrt{D_{jj}}z_j\) using independent
standard normals. Multiplication by \(D\) is a vector scaling operation;
there is no reason to construct a dense diagonal matrix.

To see why the draw is correct, let \(K=DX\trans H^{-1}\). Its mean is
\(K\widetilde y\), and the Woodbury identity gives
\begin{equation}\label{eq:woodbury}
 A^{-1}=D-DX\trans H^{-1}XD,
 \qquad A^{-1}X\trans=DX\trans H^{-1}=K.
\end{equation}
Also \(\Cov(u,v)=DX\trans\) and \(\Var(v)=H\). Consequently,
\[
 \Var(u-Kv)=D-KXD-DX\trans K\trans+KHK\trans
           =D-DX\trans H^{-1}XD=A^{-1}.
\]
Multiplication by \(\sigma\) gives precisely the mean and covariance in
\eqref{eq:beta}. This is a different way to sample the same conditional,
not an approximation to it.

For a dense design, forming \(XDX\trans\) costs \(O(n^2p)\) and its
factorization costs \(O(n^3)\). Storage is \(O(np+n^2)\), including the
design matrix. Thus ``linear in \(p\)'' assumes \(n\) is fixed; it is not
linear in both dimensions. For very large scales, cancellation in the final
update can still cause numerical trouble, so monitor solver accuracy.

\subsection{A coordinatewise alternative}

For completeness, each coefficient can also be updated alone. Let
\(r_{-j}=y-\alpha\ones_n-\sum_{k\ne j}X_k\beta_k\) and
\(a_j=X_j\trans X_j+1/(\tau^2\lambda_j^2)\). Then
\begin{equation}
 \beta_j\given\bbeta_{-j},\cdot\sim
 \Normal\left(\frac{X_j\trans r_{-j}}{a_j},\frac{\sigma^2}{a_j}\right).
\end{equation}
A sequential coordinate sweep is valid and can use inexpensive residual
updates. With correlated predictors, however, the chain may move slowly
along directions that require coordinated changes in several coefficients.
A block draw often improves mixing even when it has a larger cost per sweep.

\subsection{Numerical details that affect correctness}

\begin{itemize}
\item \textbf{Use the correct gamma convention.} To draw \(Z\sim\IG(a,b)\),
one option is \(G\sim\Ga(a,\text{rate}=1)\), then \(Z=b/G\).
This avoids ambiguities in a library's second parameter. Finite-precision
underflow or overflow still requires explicit detection.
\item \textbf{Record scales on the log scale.} The horseshoe can explore very
small and very large positive scales. Log-scale traces expose this behavior
more clearly. Stable evaluation of sums or products may be needed before
forming reciprocals.
\item \textbf{Do not silently clip.} Truncating \(\tau\), \(\lambda_j\),
or their reciprocals at arbitrary bounds changes the transition and generally
the target. If a truncated prior is intended, specify it and derive the
corresponding sampler.
\item \textbf{Distinguish a numerical failure from an ill-defined model.}
At finite positive scales, \(A\) is positive definite mathematically. A
failed Cholesky factorization can indicate overflow, lost symmetry, severe
conditioning, or a programming error. Unreported diagonal jitter is an
unreported change to the draw.
\item \textbf{Keep the linear algebra auditable.} Reuse fixed cross-products,
check relative solve residuals, and record which Gaussian sampler is used.
Avoid explicit inverses and avoid rounding coefficients during the chain.
\end{itemize}

\section{Choosing prior scales and checking their meaning}\label{sec:hyperparameters}

\subsection{Calibrating global shrinkage to anticipated sparsity}

Let \(p_0\) be a plausible number of substantial effects, with \(0<p_0<p\).
Under an orthogonal design with common squared column norm \(c\), define
the effective number of unshrunk coordinates as
\begin{equation}\label{eq:meff}
 m_{\rm eff}=\sum_{j=1}^p(1-\kappa_j).
\end{equation}
This is a continuous summary in \((0,p)\), not an integer model size. The
sparsity calibration idea follows Piironen and Vehtari \cite{piironen2017};
the algebra below uses our noise-scaled convention.

At fixed \(\tau\), put \(t=\sqrt c\tau\). Since
\(\lambda\sim\HC(0,1)\),
\begin{align}
 \E(\kappa\given\tau)
 &=\frac2\pi\int_0^\infty
       \frac{d\lambda}{(1+\lambda^2)(1+t^2\lambda^2)}
   =\frac{1}{1+t},\\
 \E(m_{\rm eff}\given\tau)&=p\frac{\sqrt c\tau}{1+\sqrt c\tau}.
\end{align}
Solving \(\E(m_{\rm eff}\given\tau_0)=p_0\) yields
\begin{equation}\label{eq:tau0}
 \tau_0=\frac{p_0}{p-p_0}\frac1{\sqrt c}.
\end{equation}
For columns with \(X_j\trans X_j=n\), a useful calibration is
\begin{equation}\label{eq:stau}
 s_\tau=\frac{p_0}{p-p_0}\frac1{\sqrt n},
 \qquad \tau\sim\HC(0,s_\tau).
\end{equation}
For example, \(n=100,p=50,p_0=5\) gives \(s_\tau=1/90\).
This is an illustrative choice of prior scale, not an empirical estimate.

Equation \eqref{eq:tau0} equates an expectation with \emph{fixed}
\(\tau=\tau_0\). Giving \(\tau\) a half-Cauchy prior with scale
\(\tau_0\) does not make the fully averaged prior mean of \(m_{\rm eff}\)
equal to \(p_0\). It sets a characteristic global scale (the half-Cauchy
median), while retaining considerable uncertainty. Correlated designs and
unequal column norms also make this a calibration heuristic rather than an
exact model-size statement.

\subsection{Why some references include a factor of \texorpdfstring{\(\sigma\)}{sigma}}

An alternative writes
\(\beta_j\given\tau_{\rm abs},\lambda_j\sim
\Normal(0,\tau_{\rm abs}^2\lambda_j^2)\), with no explicit \(\sigma^2\)
in the prior variance. Then the orthogonal shrinkage factor is
\[
 \kappa_j=\frac{1}{1+c_j\tau_{\rm abs}^2\lambda_j^2/\sigma^2},
\]
and the analogous reference scale is
\(\tau_{{\rm abs},0}=\sigma p_0/\{(p-p_0)\sqrt n\}\).
In our notation \(\tau_{\rm abs}=\sigma\tau\), so that factor of
\(\sigma\) is already accounted for. Multiplying \eqref{eq:stau} by
\(\sigma\) again would double-count it. An independent prior on
\(\tau_{\rm abs}\) is also a different joint model from a prior on
\(\tau\) independent of \(\sigma\).

\subsection{The noise-variance prior}

The main derivation uses positive \(a_\sigma,b_\sigma\). This makes the
variance prior proper and keeps its units explicit: \(b_\sigma\) has units
of response squared. If \(a_\sigma>1\), its mean is
\(b_\sigma/(a_\sigma-1)\); if \(a_\sigma>2\), its variance is
\(b_\sigma^2/\{(a_\sigma-1)^2(a_\sigma-2)\}\).

Choose parameters using plausible noise levels on the analysis scale, and
check the implied prior distribution. For example,
\(\IG(2,s_0^2)\) has mean \(s_0^2\) but infinite variance; it is a specific
heavy-tailed prior, not an automatically innocuous default. Very small
inverse-gamma shape and scale values should not be called ``noninformative''
without examining their implications.

The common scale-invariant prior \(p(\sigma^2)\propto1/\sigma^2\) is
improper. Its conditional formula is obtained formally by setting
\(a_\sigma=b_\sigma=0\), but this is not a proper \(\IG(0,0)\)
distribution. Properness of the resulting posterior must be established for
the actual model and data; valid-looking conditional distributions alone do
not prove it. A half-Cauchy prior on \(\sigma\) is another conjugately
augmentable option, derived in Appendix \ref{app:variants}.

\subsection{Prior predictive reasoning and sensitivity}

The prior predicts a regression surface \(\alpha\ones_n+X\bbeta\) and
responses around it. Ask whether the implied effect sizes and response ranges
are scientifically plausible. A literal prior predictive distribution
requires a proper intercept prior or a fixed baseline; the flat intercept
prior in the main derivation cannot itself generate a proper prior predictive
distribution.

For a future fitted analysis, vary \(p_0\) or \(s_\tau\), the noise prior,
and scientifically relevant predictor scaling choices. Compare substantive
coefficient summaries, selected sets, and predictive performance. Sensitivity
is especially important when the data weakly identify the global scale or
when many predictors are strongly correlated.

\section{What to compute from posterior draws}\label{sec:inference}

Let \(s=1,\ldots,S\) index retained draws from adequately explored chains.
For a posterior functional \(g\), a Monte Carlo approximation is
\begin{equation}
 \E[g(\theta)\given y,X]\ \approx\ \frac1S\sum_{s=1}^Sg(\theta^{(s)}),
\end{equation}
provided the relevant expectation exists. Dependence between draws affects
Monte Carlo precision, which is measured through effective sample sizes and
Monte Carlo standard errors rather than through \(S\) alone.

\subsection{Coefficient summaries and uncertainty}

Useful summaries include posterior medians, central credible intervals,
probabilities of a positive effect, and probabilities of exceeding a
scientifically meaningful magnitude. A posterior mean is useful under squared
error loss when it exists. In weakly identified directions under heavy-tailed
priors, moment existence and numerical stability should not be assumed from
the fact that a finite Monte Carlo average can be calculated.

A central \(100(1-\gamma)\%\) interval for \(\beta_j\) is the interval
between its empirical \(\gamma/2\) and \(1-\gamma/2\) posterior quantiles.
It describes uncertainty conditional on the model, data, and prior. It is not
automatically a simultaneous statement about all \(p\) coefficients, and
does not automatically have nominal repeated-sampling coverage.

When using the orthogonal interpretation, summarize
\(\kappa_j^{(s)}=1/\{1+c_j\tau^{2(s)}\lambda_j^{2(s)}\}\) and
\(m_{\rm eff}^{(s)}=\sum_j(1-\kappa_j^{(s)})\). In general,
\[
 \E\left[\frac1{1+c_j\tau^2\lambda_j^2}\given y\right]
 \ne \frac1{1+c_j\E(\tau^2\given y)\E(\lambda_j^2\given y)}.
\]
The right-hand side also relies on moments that may be unstable or absent.
Compute the nonlinear quantity within each draw and then summarize it.

\subsection{Posterior prediction}

Transform a new predictor vector with the recorded training-data constants.
For every posterior draw, compute
\begin{equation}
 \mu_*^{(s)}=\alpha^{(s)}+x_*\trans\bbeta^{(s)},\qquad
 y_*^{(s)}=\mu_*^{(s)}+\sigma^{(s)}z_*^{(s)},\quad
 z_*^{(s)}\sim\Normal(0,1).
\end{equation}
The distribution of \(\mu_*^{(s)}\) describes uncertainty about the conditional
mean. The distribution of \(y_*^{(s)}\) additionally includes observation
noise. For multiple new observations, use the same parameter draw for the
whole vector and independent noise draws conditional on that state; parameter
uncertainty then induces dependence between their predictive distributions.

When the relevant moments exist, total variance gives
\begin{equation}
 \Var(y_*\given y,X)=\E(\sigma^2\given y,X)
       +\Var(\alpha+x_*\trans\bbeta\given y,X).
\end{equation}
This identity explains why plugging in coefficients and adding only a noise
interval understates predictive uncertainty.

\subsection{A general-design measure of conditional complexity}

For the penalized slopes, conditional fitted values use the matrix
\(S_D=X(X\trans X+D^{-1})^{-1}X\trans\). Its trace
\begin{equation}
 d_{\rm eff}(D)=\tr(S_D)
\end{equation}
is a conditional smoothing-complexity summary. In the orthogonal case it is
\(\sum_j(1-\kappa_j)\), but in a correlated design the trace accounts for
joint directions in predictor space. It is not an exact count of selected
variables and is not by itself the degrees of freedom of the fully adaptive
posterior procedure. For centered predictors, an unpenalized estimated
intercept contributes an additional one to the conditional smoother trace.

\section{Variable selection is a posterior decision}\label{sec:selection}

\subsection{Why there is no ordinary inclusion probability}

With a proper posterior under the continuous horseshoe prior,
\begin{equation}
 \Prob(\beta_j=0\given y,X)=0,\qquad
 \Prob(\beta_j\ne0\given y,X)=1.
\end{equation}
Neither probability ranks variables. A spike-and-slab model instead places
a point mass at zero and can define a latent inclusion indicator. The
ordinary horseshoe has no such indicator. Its proximity to sparse estimates
does not turn \(1-\kappa_j\) into a posterior inclusion probability.

A selection report should therefore state the event being judged, the
threshold or loss, and whether the goal is scientific interpretation or a
smaller predictive model.

\subsection{Rule 1: exclusion of zero from a credible interval}

One simple rule retains \(j\) if a chosen central credible interval for
\(\beta_j\) excludes zero. Equivalently, for a continuous posterior, a
central \(100(1-\gamma)\%\) interval excludes zero when
\[
 \Prob(\beta_j>0\given y)>1-\gamma/2
 \quad\text{or}\quad
 \Prob(\beta_j<0\given y)>1-\gamma/2.
\]
This rule rewards sign certainty. It can omit important but highly correlated
predictors whose individual effects are uncertain. It also does not control
the familywise error rate or the frequentist false discovery rate merely
because each interval has a familiar nominal level.

\subsection{Rule 2: practical relevance and an explicit loss}

Choose a threshold \(\delta_j>0\) in meaningful coefficient units and define
\begin{equation}
 q_j=\Prob(|\beta_j|>\delta_j\given y,X)
 \approx\frac1S\sum_s\ind\{|\beta_j^{(s)}|>\delta_j\}.
\end{equation}
Unlike exact nonzero probability, this is informative under a continuous
prior. For standardized predictors, \(\delta_j\) can represent a minimum
response change for a one-standard-deviation predictor change.

Let \(d_j=1\) mean selecting predictor \(j\). Suppose a false selection
costs \(C_{\rm FP}>0\), and missing a practically relevant effect costs
\(C_{\rm FN}>0\). Under additive loss, the two posterior expected losses are
\[
 L(d_j=1\given y)=C_{\rm FP}(1-q_j),\qquad
 L(d_j=0\given y)=C_{\rm FN}q_j.
\]
Select if and only if
\begin{equation}\label{eq:selectionloss}
 q_j>\frac{C_{\rm FP}}{C_{\rm FP}+C_{\rm FN}},
\end{equation}
with a declared convention at equality. This threshold follows from the
specified loss, rather than from an arbitrary label of ``included.''

For a nonempty reported set \(S_0\), the posterior expected fraction of
practically negligible selected effects is
\begin{equation}
 \E\left[\frac{\sum_{j\in S_0}\ind\{|\beta_j|\leq\delta_j\}}
                  {|S_0|}\given y,X\right]
 =\frac1{|S_0|}\sum_{j\in S_0}(1-q_j).
\end{equation}
One may sort \(q_j\)'s and retain a largest prefix whose average
\(1-q_j\) does not exceed a chosen tolerance. Linearity of expectation makes
this formula valid without posterior independence. It is a Bayesian statement
about practical-null events under the model, not an automatic frequentist
false discovery guarantee.

\subsection{Rule 3: a shrinkage-score threshold}

An exploratory rule ranks variables by
\(r_j=\E(1-\kappa_j\given y,X)\), possibly selecting those above a stated
cutoff. This is easy to compute but depends on column scaling and on the
orthogonal interpretation. A cutoff such as \(1/2\) is a convention unless
justified by an explicit decision problem. Report it as a shrinkage-score
rule, not as a posterior probability that the variable belongs to a model.

\subsection{A smaller predictive model and the effect of correlation}

If the goal is a compact predictor, a useful target is to approximate the
full model's fitted or predictive distribution using fewer variables. For
example, for a candidate subset \(S_0\) and each posterior draw, one can
minimize
\[
 \min_{a,b_{S_0}}\frac1n
 \norm{\alpha^{(s)}\ones_n+X\bbeta^{(s)}
                -a\ones_n-X_{S_0}b_{S_0}}^2.
\]
This defines an empirical squared-error projection of the full fitted mean;
it is not a complete predictive selection procedure by itself. Choosing a
subset size still requires predictive assessment, and uncertainty or residual
variance for the reduced predictor must account for approximation error.
Perform subset search and tuning within training folds before evaluating on
held-out data.

In a correlated group, retaining one representative may preserve predictions
even when the data cannot distinguish which member has an individual effect.
For scientific interpretation, report that ambiguity, consider prespecified
group contributions, and avoid interpreting a selected slope as a causal
effect without additional assumptions.

\key{What a transparent selection result should contain}
State the preprocessing, prior, decision rule, interval level or practical
threshold, selected set, and uncertainty summaries. Include sensitivity to
reasonable thresholds and explain how correlated predictors were handled.
If a reduced model is refitted, describe its new inferential target rather
than treating the refit as the original horseshoe posterior.

\section{Diagnostics and a validation plan}\label{sec:diagnostics}

\subsection{Convergence and Monte Carlo accuracy}

Use several independently seeded chains with dispersed reasonable starts.
Inspect traces and between-chain agreement for coefficients, \(\log\tau\),
\(\log\sigma\), representative local scales, fitted values, and important
decisions. A coefficient near zero can look stable while its global and
local scales remain highly dependent and explore slowly.

Rank-normalized split \(\widehat R\), bulk effective sample size, and tail
effective sample size are useful diagnostics for heavy-tailed posteriors
\cite{vehtari2021}. Values of \(\widehat R\) close to one, often with a
working target below \(1.01\), are screening criteria rather than proofs of
convergence. Assess effective sample sizes and Monte Carlo errors for the
actual reported quantities, including interval endpoints and selection-event
probabilities. A threshold decision is unreliable if its Monte Carlo error
is comparable to its distance from the threshold.

Discard an initial transient based on exploration and diagnostics, and
extend runs when necessary. Routine thinning discards information; it is
mainly a storage choice, not a remedy for slow mixing. Prefer improving
parameterization or block updates and keeping enough draws for accurate
summaries. There is no single iteration count that guarantees adequacy for
every design and signal configuration.

\subsection{Posterior predictive checks}

From retained states generate replicated observations
\(y_i^{\rm rep,(s)}\sim
\Normal(\alpha^{(s)}+x_i\trans\bbeta^{(s)},\sigma^{2(s)})\).
Compare observed and replicated summaries that address the likelihood:
residual spread versus fitted value, extreme residuals, skewness, and
dependence across time or groups where relevant. A well-mixed chain can still
fit a misspecified homoscedastic Gaussian model. Predictive checking and
sampler convergence answer different questions.

\subsection{Checks for a future implementation}

The following are proposed checks, not completed experiments in these notes:
\begin{enumerate}
\item \textbf{Distribution convention.} At fixed \(a,b\), confirm that the
inverse-gamma generator agrees with the reciprocal-gamma relationship, using
quantiles rather than nonexistent moments when \(a\leq1\).
\item \textbf{Gaussian block at fixed scales.} Compare sampled means and
covariances against the exact normal target for a small problem. Use a
non-diagonal precision so a transposed-Cholesky bug is detectable. Compare
both Gaussian algorithms within Monte Carlo error.
\item \textbf{Orthogonal reduction.} Check that the block sampler reproduces
\eqref{eq:orthposterior} for orthogonal columns, with the correct column
norm. Include \(c_j=n\), not only \(c_j=1\).
\item \textbf{Conditional log densities.} Vary one block in
\eqref{eq:jointkernel} while holding the others fixed and verify agreement,
up to an additive constant, with its proposed full-conditional log density.
This catches omitted normalizing factors.
\item \textbf{Intercept and rank cases.} Compare the explicit and correctly
integrated intercept formulations; include correlated, duplicated, and
\(p>n\) designs. Check both predictions and noise-variance behavior.
\item \textbf{End-to-end calibration.} With proper priors for every parameter,
including the intercept, consider simulation-based calibration and comparisons
with a separately implemented posterior method. Such tests must account for
Monte Carlo dependence and uncertainty.
\end{enumerate}

\subsection{Reproducibility for future computational work}

Any future simulations, benchmarks, or computational validation for this
project should run as persistent jobs on \texttt{owls.stat.rice.edu}, after
reading its applicable \texttt{AGENTS.md} instructions. Local editing and
document compilation do not constitute a sampler experiment. These notes
contain no claim that a numerical sampler has been run.

A computational run should preserve the exact source and configuration,
data/preprocessing description, random seeds and generator details, SHA-256
hashes, launch command, scheduler or process identifiers, environment and
package versions, timestamps, exit status, logs, generated data, retained
draws, and diagnostic evidence. Conclusions should be based on inspected
saved outputs. If the remote system is unavailable, report that fact instead
of substituting a local experiment.

\appendix
\clearpage
\section{Intercept integration and collapsed updates}\label{app:intercept}

\subsection{Integrating out the flat intercept exactly}

Assume \(n\geq2\), \(\ones_n\trans X=0\), and set
\(y_c=y-\bar y\ones_n\). The residual identity is
\[
 \norm{y-\alpha\ones_n-X\bbeta}^2
 =\norm{y_c-X\bbeta}^2+n(\alpha-\bar y)^2.
\]
Integrating the likelihood over \(\alpha\) introduces the factor
\(\sqrt{2\pi\sigma^2/n}\). Therefore the integrated likelihood is
\begin{equation}
 p(y\given\bbeta,\sigma^2)\ \propto\
 (\sigma^2)^{-(n-1)/2}
 \exp\left\{-\frac{\norm{y_c-X\bbeta}^2}{2\sigma^2}\right\},
\end{equation}
where the proportionality reflects the arbitrary constant in the flat
intercept prior. The coefficient update uses \(\widetilde y=y_c\), while
the variance update is
\begin{equation}\label{eq:integratedalpha}
 \sigma^2\given\bbeta,D,y\sim
 \IG\left(a_\sigma+\frac{n-1+p}{2},\ b_\sigma+
           \frac{\norm{y_c-X\bbeta}^2+\bbeta\trans D^{-1}\bbeta}{2}\right).
\end{equation}
Other scale updates are unchanged. To recover intercept uncertainty, draw
\(\alpha\sim\Normal(\bar y,\sigma^2/n)\) conditional on retained states.
For uncentered predictors the analogous formula uses the centering projector
\(M=I_n-\ones_n\ones_n\trans/n\), so replace \(X\) and \(y\) by
\(MX\) and \(My\).

Subtracting \(\bar y\) from the response but using an \(n\)-dimensional
zero-intercept likelihood with the original \(n/2\) power of \(\sigma^2\)
would instead treat the estimated baseline as known. That is a different
variance calculation.

\subsection{Integrating out the coefficients for a joint block}

Hold \(\alpha,D\) fixed and integrate out \(\bbeta\). Normal convolution
gives
\begin{equation}
 \widetilde y\given\sigma^2,D\sim\Normal_n(0,\sigma^2H),
 \qquad H=I_n+XDX\trans.
\end{equation}
Therefore
\begin{equation}\label{eq:collapsedbeta}
 \sigma^2\given\alpha,D,y\sim
 \IG\left(a_\sigma+\frac n2,\ b_\sigma+
          \frac{\widetilde y\trans H^{-1}\widetilde y}{2}\right).
\end{equation}
Here there is no \(p/2\): the coefficients have actually been integrated
out, including their normalizing factors. This explains how two variance
updates with different shapes can both be correct for different conditionals.

A valid joint-block replacement for the main sweep is:
\begin{enumerate}
\item update \(\alpha\) conditional on the current \(\bbeta,\sigma^2\);
\item draw \(\sigma^2\) using \eqref{eq:collapsedbeta}, conditional on the
      current \(\alpha,D\);
\item immediately draw \(\bbeta\) using \eqref{eq:beta} with that new
      \(\sigma^2\), before updating any scale conditioned on \(\bbeta\);
\item update the local, global, and auxiliary scales as before.
\end{enumerate}
Steps 2--3 together draw from the joint conditional of
\((\sigma^2,\bbeta)\). Inserting a collapsed variance formula somewhere
inside an otherwise unchanged scan without respecting the relevant block
ordering can produce an invalid chain.

\subsection{Properness under the main variance prior}

With \(n\geq1\), \(a_\sigma,b_\sigma>0\), proper scale priors, and the
flat intercept prior, the posterior in the main model is proper. One way to
see this is to integrate \(\alpha\) first. Its integrated likelihood is at
most a constant times \((\sigma^2)^{-(n-1)/2}\), since the residual
exponential is at most one. Integrating the proper coefficient and scale
priors cannot increase that bound. Finally,
\[
 \int_0^\infty
 (\sigma^2)^{-(n-1)/2}\,p(\sigma^2)\,d\sigma^2<\infty
\]
under a proper inverse-gamma prior with positive scale. The likelihood is
positive for finite data, so the normalizing constant is also positive.
This argument does not carry over automatically to
\(p(\sigma^2)\propto1/\sigma^2\), and properness alone does not ensure
existence of every posterior moment.

\section{Alternative prior choices change specific updates}\label{app:variants}

\subsection{A coefficient prior independent of the noise variance}

Consider the explicitly different hierarchy
\[
 \bbeta\given\tau_{\rm abs}^2,\blambda^2
 \sim\Normal_p(0,D_{\rm abs}),\qquad
 D_{\rm abs}=\tau_{\rm abs}^2\diag(\lambda_j^2),
\]
with \(\tau_{\rm abs}\sim\HC(0,s_{\rm abs})\) independent of \(\sigma\).
The likelihood, flat intercept prior, and proper inverse-gamma noise prior
remain as in the main model. Using the analogous augmentation, the updated
conditionals are
\begin{align}
 V&=(X\trans X/\sigma^2+D_{\rm abs}^{-1})^{-1},\nonumber\\
 \bbeta\given\cdot&\sim\Normal_p\left(
          VX\trans\widetilde y/\sigma^2,V\right),\\
 \sigma^2\given\cdot&\sim\IG\left(a_\sigma+\frac n2,
                         b_\sigma+\frac R2\right),\\
 \lambda_j^2\given\cdot&\sim\IG\left(1,
               \frac1{\nu_j}+\frac{\beta_j^2}{2\tau_{\rm abs}^2}\right),\\
 \tau_{\rm abs}^2\given\cdot&\sim\IG\left(\frac{p+1}{2},
               \frac1\xi+\frac12\sum_j\frac{\beta_j^2}{\lambda_j^2}\right),\\
 \nu_j\given\cdot&\sim\IG(1,1+1/\lambda_j^2),\\
 \xi\given\cdot&\sim\IG(1,1/s_{\rm abs}^2+1/\tau_{\rm abs}^2).
\end{align}
The intercept update is unchanged. Relative to the main model, the noise
update loses both the coefficient-prior penalty and the \(p/2\) shape
contribution; the local and global updates lose their \(\sigma^2\) factors.
One cannot mix this variance update with the main model's coefficient update.

\subsection{A half-Cauchy prior on the noise standard deviation}

Return to the main, noise-scaled coefficient prior, but replace
\eqref{eq:otherpriors}'s variance prior by
\[
 \sigma\sim\HC(0,s_\sigma),\qquad
 \sigma^2\given\eta\sim\IG(1/2,1/\eta),\qquad
 \eta\sim\IG(1/2,1/s_\sigma^2).
\]
The resulting updates are
\begin{align}
 \sigma^2\given\cdot&\sim
 \IG\left(\frac{n+p+1}{2},\frac1\eta+\frac{R+Q}{2}\right),\\
 \eta\given\cdot&\sim\IG\left(1,\frac1{s_\sigma^2}+\frac1{\sigma^2}\right).
\end{align}
All other main-model conditionals are unchanged. These formulas condition on
the explicit intercept. The appropriate residual-dimension change is needed
if the intercept is integrated out. Because the variance prior now behaves
differently near zero, the inverse-gamma properness bound in Appendix
\ref{app:intercept} should not be applied without checking its assumptions.

\subsection{Fixing the global scale}

If \(\tau=\tau_0>0\) is fixed by the model, omit both the \(\tau^2\) and
\(\xi\) updates. Keep the remaining conditionals with that fixed value.
This is a useful mathematically simpler model, but it excludes posterior
uncertainty about global shrinkage. Choosing \(\tau_0\) by an optimization
or empirical-Bayes step is a separate procedure and should be reported.

\subsection{The regularized horseshoe}

The regularized horseshoe adds control over the largest coefficient scales
\cite{piironen2017}. In a noise-scaled convention, one possible form is
\begin{equation}
 \widetilde\lambda_j^2
  =\frac{c^2\lambda_j^2}{c^2+\tau^2\lambda_j^2},\qquad
 \beta_j\given\cdot\sim
 \Normal(0,\sigma^2\tau^2\widetilde\lambda_j^2),\qquad c>0.
\end{equation}
For \(\tau^2\lambda_j^2\ll c^2\), this approximately recovers the ordinary
horseshoe variance. For \(\tau^2\lambda_j^2\gg c^2\), the variance tends
to \(\sigma^2c^2\), limiting the scale of very large effects. Here \(c\)
is measured relative to the noise standard deviation; a slab width in
absolute response units requires a corresponding reparameterization.

This modification changes the prior normalization as a function of
\(\lambda_j^2\) and \(\tau^2\). The simple inverse-gamma local and global
updates derived for the ordinary horseshoe no longer apply unchanged.
Re-derive the target and use an appropriate sampler; substituting
\(\widetilde\lambda_j^2\) in one line of the ordinary Gibbs sweep is not
a valid derivation.

\clearpage
\section{A one-predictor conditional example}\label{app:example}

To make each quantity concrete, suppose one centered predictor has
\(X\trans X=20\), and the currently intercept-adjusted response satisfies
\(X\trans\widetilde y=10\) and \(\widetilde y\trans\widetilde y=15\).
Take a current state \(\sigma^2=2,\tau^2=0.25,\lambda^2=4\).
Then \(D=1\), \(A=21\), and the coefficient update is
\[
 \beta\given\cdot\sim\Normal(10/21,\,2/21).
\]
The unregularized coordinate estimate is \(10/20=0.5\); its conditional
shrinkage factor is \(\kappa=1/21\), so its conditional posterior mean
is \((20/21)0.5=10/21\), agreeing with the matrix calculation.

For illustration of the next conditional, evaluate the formulas at
\(\beta=0.4\), with \(n=20,p=1,a_\sigma=2,b_\sigma=1\). The residual
sum of squares and prior penalty are
\[
 R=15-2(0.4)(10)+(0.4)^2(20)=10.2,\qquad Q=(0.4)^2=0.16.
\]
Thus
\[
 \sigma^2\given\cdot\sim\IG\left(12.5,\,6.18\right).
\]
If evaluating the local-scale conditional at \(\nu=1\) and \(\sigma^2=2\),
while holding \(\tau^2=0.25\), then
\[
 \lambda^2\given\cdot\sim\IG\left(1,\,
 1+\frac{0.16}{2(2)(0.25)}\right)=\IG(1,1.16).
\]
These are separate conditional evaluations at stated values, not successive
realized Gibbs draws. In an actual sweep, the local-scale update must use
the newly drawn noise variance. No random experiment is being reported here.

\clearpage
\section{Quick reference for the main sampler}\label{app:quick}

\textbf{Model.} Explicit flat-prior intercept; noise-scaled coefficient prior;
\(\sigma^2\sim\IG(a_\sigma,b_\sigma)\);
\(\lambda_j\sim\HC(0,1)\); \(\tau\sim\HC(0,s_\tau)\).
The inverse-gamma density is proportional to \(z^{-a-1}e^{-b/z}\).

\textbf{Definitions.}
\begin{align*}
 D&=\tau^2\diag(\lambda_1^2,\ldots,\lambda_p^2),
 &A&=X\trans X+D^{-1},\\
 \widetilde y&=y-\alpha\ones_n,
 &m&=A^{-1}X\trans\widetilde y,\\
 R&=\norm{\widetilde y-X\bbeta}^2,
 &Q&=\sum_j\frac{\beta_j^2}{\tau^2\lambda_j^2}.
\end{align*}

\textbf{Full conditionals.} In a sequential scan, always use current values.
\begin{align*}
 \alpha\given\cdot&\sim\Normal\left(
        \frac{\ones_n\trans(y-X\bbeta)}n,\frac{\sigma^2}n\right),\\[4pt]
 \bbeta\given\cdot&\sim\Normal_p(m,\sigma^2A^{-1}),\\[4pt]
 \sigma^2\given\cdot&\sim\IG\left(a_\sigma+\frac{n+p}{2},
                               b_\sigma+\frac{R+Q}{2}\right),\\[4pt]
 \lambda_j^2\given\cdot&\sim\IG\left(1,
                    \frac1{\nu_j}+\frac{\beta_j^2}{2\sigma^2\tau^2}\right),\\[4pt]
 \tau^2\given\cdot&\sim\IG\left(\frac{p+1}{2},
                \frac1\xi+\frac1{2\sigma^2}\sum_j\frac{\beta_j^2}{\lambda_j^2}\right),\\[4pt]
 \nu_j\given\cdot&\sim\IG\left(1,1+\frac1{\lambda_j^2}\right),\\[4pt]
 \xi\given\cdot&\sim\IG\left(1,\frac1{s_\tau^2}+\frac1{\tau^2}\right).
\end{align*}

\textbf{Interpretation.} Small \(\tau\) encourages global shrinkage; large
\(\lambda_j\) allows a local exception. An orthogonal-design shrinkage factor
is \(\kappa_j=1/\{1+(X_j\trans X_j)\tau^2\lambda_j^2\}\).
The horseshoe is continuous; exact-zero selection requires a separate
decision rule or a different prior.

\textbf{Most consequential mistakes.} Confusing gamma rate and scale;
omitting the \(p/2\) or \(Q\) term in the main noise update; shrinking the
intercept unintentionally; forgetting column norms in \(\kappa_j\);
using the wrong Cholesky transpose; replacing random draws by means;
simultaneously updating dependent blocks with stale states; and reporting a
shrinkage score as an inclusion probability.

\clearpage
\phantomsection
\addcontentsline{toc}{section}{References}
\begin{thebibliography}{9}
\raggedright
\bibitem{carvalho2009}
Carvalho, C. M., Polson, N. G., and Scott, J. G. (2009).
Handling sparsity via the horseshoe.
\emph{Proceedings of the Twelfth International Conference on Artificial
Intelligence and Statistics}, PMLR \textbf{5}, 73--80.
\href{https://proceedings.mlr.press/v5/carvalho09a.html}{Paper and open-access PDF}.

\bibitem{carvalho2010}
Carvalho, C. M., Polson, N. G., and Scott, J. G. (2010).
The horseshoe estimator for sparse signals.
\emph{Biometrika}, \textbf{97}(2), 465--480.
\href{https://doi.org/10.1093/biomet/asq017}{doi:10.1093/biomet/asq017}.

\bibitem{makalic2016}
Makalic, E. and Schmidt, D. F. (2016).
A simple sampler for the horseshoe estimator.
\emph{IEEE Signal Processing Letters}, \textbf{23}(1), 179--182.
\href{https://doi.org/10.1109/LSP.2015.2503725}{doi:10.1109/LSP.2015.2503725}.
\href{https://arxiv.org/abs/1508.03884}{Open-access preprint}.

\bibitem{bhattacharya2016}
Bhattacharya, A., Chakraborty, A., and Mallick, B. K. (2016).
Fast sampling with Gaussian scale-mixture priors in high-dimensional regression.
\emph{Biometrika}, \textbf{103}(4), 985--991.
\href{https://doi.org/10.1093/biomet/asw042}{doi:10.1093/biomet/asw042}.
\href{https://arxiv.org/abs/1506.04778}{Open-access preprint}.

\bibitem{piironen2017}
Piironen, J. and Vehtari, A. (2017).
Sparsity information and regularization in the horseshoe and other shrinkage priors.
\emph{Electronic Journal of Statistics}, \textbf{11}(2), 5018--5051.
\href{https://doi.org/10.1214/17-EJS1337SI}{doi:10.1214/17-EJS1337SI}.
\href{https://arxiv.org/abs/1707.01694}{Open-access preprint}.

\bibitem{vehtari2021}
Vehtari, A., Gelman, A., Simpson, D., Carpenter, B., and B\"urkner, P.-C. (2021).
Rank-normalization, folding, and localization: An improved
\(\widehat R\) for assessing convergence of MCMC (with discussion).
\emph{Bayesian Analysis}, \textbf{16}(2), 667--718.
\href{https://doi.org/10.1214/20-BA1221}{doi:10.1214/20-BA1221}.
\href{https://arxiv.org/abs/1903.08008}{Open-access preprint}.
\end{thebibliography}
\end{document}
